% 本文件由 LaTeX 排版 Agent 根据有效 translation.json 自动生成。
% author=Hippolytus; work_id=0524; title=论宗徒与门徒

\chapter{\WorkTitle{论宗徒与门徒}}

% structure: p0001 source=h1 target=omitted reason=page-title-already-rendered-by-parent

% structure: p0002 source=h2 target=section reason=relative-heading-rank; base=h2

\section{\Emph{论十二宗徒}}

他们每人在何处宣讲，以及何处是他的结局。\par

1. \Person{伯多禄}在本都、迦拉达、卡帕多细雅、贝塔尼亚、意大利和亚细亚宣讲福音，后来在罗马被\Person{尼禄}倒钉十字架，正如他本人所愿的那样受苦。\par

2. \Person{安德肋}向斯基泰人和色雷斯人宣讲，在阿哈雅的\Emph{一个城镇}帕特雷被钉在橄榄树上而死，并埋葬在那里。\par

3. \Person{若望}在亚细亚被\Person{多米提安}王放逐到帕特摩岛，在那里写了福音并看见了默示录；在图拉真时代，他在厄弗所安眠，人们寻找他的遗骸却未能找到。\par

4. \Person{雅各伯}，他的兄弟，在犹太宣讲时，被分封侯\Person{黑落德}用刀处死，并埋葬在那里。\par

5. \Person{斐理伯}在夫黎基雅宣讲，在\Person{多米提安}时代于希拉波利斯被倒钉十字架而死，并埋葬在那里。\par

6. \Person{巴尔多禄茂}向印度人宣讲，也将\BibleBook{玛窦}{福音}给了他们，被倒钉十字架，埋葬在大亚美尼亚的\Emph{一个城镇}阿拉努姆。\par

7. \Person{玛窦}用希伯来语写了福音，在耶路撒冷出版，并在帕提雅的\Emph{一个城镇}希埃雷斯安眠。\par

8. \Person{多默}向帕提雅人、玛待人、波斯人、希尔卡尼亚人、巴克特里亚人和马尔吉亚那人宣讲，在印度的\Place{卡拉梅内}城被松木矛刺穿四肢而死，并埋葬在那里。\par

9. \Person{阿耳斐的儿子雅各伯}在耶路撒冷宣讲时，被犹太人用石头打死，并埋葬在圣殿附近。\par

10. \Person{犹达}，也\Emph{称为}\Person{勒拜}，向厄得撒人和整个美索不达米亚宣讲，在贝鲁特安眠，并埋葬在那里。\par

11. \Person{热诚者西满}，\Person{克罗帕}的儿子，也\Emph{称为}\Person{犹达}，在\Person{义人雅各伯}之后成为耶路撒冷的主教，在那里安眠并被埋葬，享年一百二十岁。\par

12. \Person{玛弟亚}原是七十人之一，与十一宗徒一同被列数，在耶路撒冷宣讲，在那里安眠并被埋葬。\par

13. \Person{保禄}在基督升天后一年加入宗徒行列，从耶路撒冷开始，一直进至依里利孔、意大利和西班牙，宣讲福音三十五年。在\Person{尼禄}时代，他在罗马被斩首，并埋葬在那里。\par

% structure: p0017 source=h2 target=section reason=relative-heading-rank; base=h2

\section{\Emph{论七十门徒}}

1. \Person{主的兄弟雅各伯}，耶路撒冷主教。\par

2. \Person{克罗帕}，耶路撒冷主教。\par

3. \Person{玛弟亚}，补充了十二宗徒中的空缺位置。\par

4. \Person{达陡}，将书信送给\Place{奥加鲁斯}。\par

5. \Person{阿纳尼雅}，为\Person{保禄}施洗，\Emph{并}成为大马士革主教。\par

6. \Person{斯德望}，首位殉道者。\par

7. \Person{斐理伯}，为太监施洗。\par

8. \Person{普洛苛洛}，尼科美底亚主教，也是第一个与女儿们一同相信而去世的人。\par

9. \Person{尼加诺尔}在\Person{斯德望}殉道时去世。\par

10. \Person{提摩}，波斯特拉主教。\par

11. \Person{帕尔默纳}，索利主教。\par

12. \Person{尼苛劳}，撒玛黎雅主教。\par

13. \Person{巴尔纳伯}，米兰主教。\par

14. \Person{马尔谷}，福音作者，亚历山大主教。\par

15. \Person{路加}，福音作者。\par

这两人属于七十门徒，因基督所说的这句话而分散：\Quote{除非人吃我的肉，喝我的血，他不配于我。}但一人因\Person{伯多禄}的引导而转向主，另一人因\Person{保禄}的引导，他们得蒙尊荣宣讲那福音，为此他们也殉道了，一人被烧死，另一人被钉在橄榄树上。\par

16. \Person{息拉}，格林多主教。\par

17. \Person{息耳瓦诺}，得撒洛尼主教。\par

18. \Person{克雷斯森斯}（\OriginalTerm{克雷斯森斯}{Crescens}），高卢的加尔刻顿主教。\par

19. \Person{厄派乃托}，迦太基主教。\par

20. \Person{安德洛尼科}，潘诺尼亚主教。\par

21. \Person{安普里亚}，奥德苏斯主教。\par

22. \Person{乌尔班}，马其顿主教。\par

23. \Person{斯塔基}，拜占庭主教。\par

24. \Person{巴尔纳伯}，赫拉克利亚主教。\par

25. \Person{菲革洛}，厄弗所主教。他曾属于\Person{西满}的派别。\par

26. \Person{赫摩革乃}。他也与前一人同心。\par

27. \Person{德玛斯}，后来成为偶像的司祭。\par

28. \Person{阿培肋}，斯米纳主教。\par

29. \Person{阿黎斯托步罗}，不列颠主教。\par

30. \Person{纳尔基索}，雅典主教。\par

31. \Person{黑落狄雍}，塔尔索主教。\par

32. \Person{阿加波}，先知。\par

33. \Person{鲁富}，忒拜主教。\par

34. \Person{阿斯恩克里托}，希尔卡尼亚主教。\par

35. \Person{弗肋贡}，马拉松主教。\par

36. \Person{赫尔默}，达尔马提亚主教。\par

37. \Person{帕特罗布罗}，普特奥利主教。\par

38. \Person{赫尔马斯}，斐理伯主教。\par

39. \Person{理诺}，罗马主教。\par

40. \Person{加约}，厄弗所主教。\par

41. \Person{斐罗洛各}，西诺普主教。\par

42, 43. \Person{奥林普}和\Person{罗狄雍}在罗马殉道。\par

44. \Person{路爵}，叙利亚的劳狄刻雅主教。\par

45. \Person{雅松}，塔尔索主教。\par

46. \Person{索息帕特}，依科尼雍主教。\par

47. \Person{特尔爵}，依科尼雍主教。\par

48. \Person{厄辣斯托}，帕内拉斯主教。\par

49. \Person{夸尔托}，贝鲁特主教。\par

50. \Person{阿颇罗}，凯撒勒雅主教。\par

51. \Person{刻法}。\par

52. \Person{索斯特乃}，科洛丰主教。\par

53. \Person{提希苛}，科洛丰主教。\par

54. \Person{厄帕夫洛狄托}，安德里亚克主教。\par

55. \Person{凯撒}，都拉基乌姆主教。\par

56. \Person{马尔谷}，\Person{巴尔纳伯}的表弟，阿波罗尼亚主教。\par

57. \Person{犹斯托}，厄娄特洛波利斯主教。\par

58. \Person{阿尔特玛}，吕斯特辣主教。\par

59. \Person{克肋孟}，撒丁岛主教。\par

60. \Person{敖乃息佛}，科罗内主教。\par

61. \Person{提希苛}，加采东主教。\par

62. \Person{卡尔颇}，色雷斯的贝鲁特主教。\par

63. \Person{厄沃多}，安提约基雅主教。\par

64. \Person{阿黎斯塔苛}，阿帕美亚主教。\par

65. \Person{马尔谷}（也称\Person{若望}），比布洛波利斯主教。\par

66. \Person{则纳}，狄奥斯波利斯主教。\par

67. \Person{费肋孟}，迦萨主教。\par

68, 69. \Person{阿黎斯塔苛}和\Person{普得}。\par

70. \Person{特洛斐摩}，与\Person{保禄}一同殉道。\par

\SourceRecord{Translated by J.H. MacMahon.；Ante-Nicene Fathers；Vol. 5.；Edited by Alexander Roberts, James Donaldson, and A. Cleveland Coxe.；Buffalo, NY: Christian Literature Publishing Co.,；1886.；<http://www.newadvent.org/fathers/0524.htm>.}
